\section{طراحی در سطح مدار}
\label{sec:book-circuit-design}
\subsection{جریان و خازن چگونه تأخیر را تعیین می‌کنند؟}
در \lr{Circuit Design}، برای تابع منطقی مشخص، شبکه ترانزیستورها و اندازه آن‌ها
انتخاب می‌شود. خروجی یک گیت، ورودی گیت بعدی و سیم میان آن‌ها را راه‌اندازی می‌کند.
عایق \lr{Gate} موجب رفتار خازنی ورودی است. خازن‌های \lr{Diffusion} و سیم نیز به بار
خروجی افزوده می‌شوند. بنابراین صفر و یک‌شدن خروجی به جابه‌جایی بار الکتریکی نیاز دارد
\cite[بخش \secnum{۱.۹}، شکل \secnum{۱.۶۰}]{weste}.

اگر خازن معادل بار را \(C_L\) بنامیم، در مدل خازن ثابت رابطه بار و ولتاژ
\(Q=C_LV\) است. با تغییر ولتاژ به اندازه \(\Delta V\)، بار به اندازه
\(\Delta Q=C_L\Delta V\) تغییر می‌کند. جریان متوسط، بار جابه‌جاشده تقسیم بر زمان است.
از این دو تعریف، بدون نیاز به حل معادله دیفرانسیل، داریم
\begin{equation}
 I_{\mathrm{avg}}=\frac{\Delta Q}{\Delta t}
 \quad\Longrightarrow\quad
 \Delta t=\frac{C_L\Delta V}{I_{\mathrm{avg}}}.
 \label{eq:impl-charge-delay}
\end{equation}
برای جریان لحظه‌ای همین معنا به صورت \(i=C_L\,dV/dt\) نوشته می‌شود.
نماد \(dV/dt\) سرعت تغییر ولتاژ را بیان می‌کند. برای تخمین زمان تغییر کامل ولتاژ
از صفر تا \(V_{DD}\)، کتاب از تقریب
\begin{equation}
 t_{\mathrm{sw}}\approx\frac{C_LV_{DD}}{I_{\mathrm{avg}}}
 \label{eq:impl-full-swing}
\end{equation}
استفاده می‌کند. جریان ترانزیستور در طول گذار تغییر می‌کند، بنابراین انتخاب
\(I_{\mathrm{avg}}\) باید با همان بازه ولتاژ سازگار باشد. \lr{Propagation Delay}
معمولاً میان نقاط مرجع گذار ورودی و خروجی، مانند نصف ولتاژ تغذیه، اندازه‌گیری می‌شود.
رابطه \eqref{eq:impl-full-swing} یک تخمین مرتبه اول برای زمان تغییر ولتاژ است.

کنترل واحدها نتیجه را روشن می‌کند. چون \(\unit{\farad}=\unit{\coulomb\per\volt}\)
و \(\unit{\ampere}=\unit{\coulomb\per\second}\)، واحد
\(C_LV/I\) برابر ثانیه می‌شود. بار بیشتر، زمان بیشتر می‌خواهد. جریان قوی‌تر همان بار
را سریع‌تر جابه‌جا می‌کند. در مدل ایدئال، اگر \(C_L=0\) باشد این سهم از تأخیر صفر
می‌شود. اگر جریان به صفر نزدیک شود، زمان لازم افزایش می‌یابد.

در مثال فرضی، \(C_L=\qty{20}{\femto\farad}\)، \(V_{DD}=\qty{1}{\volt}\)
و \(I_{\mathrm{avg}}=\qty{100}{\micro\ampere}\) هستند. جای‌گذاری مرحله‌ای می‌دهد
\begin{equation*}
 t_{\mathrm{sw}}\approx
 \frac{20\times10^{-15}\times1}{100\times10^{-6}}\,\unit{\second}
 =2\times10^{-10}\,\unit{\second}=\qty{200}{\pico\second}.
\end{equation*}
دوبرابرشدن بار با جریان ثابت، زمان را به \qty{400}{\pico\second} می‌رساند.
این اعداد صرفاً برای تمرین تبدیل پیشوند و تحلیل رابطه انتخاب شده‌اند.

\subsection{اثر عرض ترانزیستور و بار مرحله بعد}
با ثابت‌بودن طول و شرایط بایاس، مدل مقدماتی جریان و خازن ورودی را متناسب با عرض
\(W\) می‌گیرد \cite[صفحه ۴۲]{weste}. افزایش عرض، جریان راه‌انداز و هم‌زمان بار
دیده‌شده در مرحله قبل را افزایش می‌دهد. اگر ضریب بزرگ‌کردن ترانزیستور \(s>0\)
باشد، برای تحلیل ساده می‌توان نوشت
\begin{equation}
 I(s)\approx sI_0,\qquad
 C_L(s)\approx C_{\mathrm{ext}}+sC_{\mathrm{int},0},\qquad
 t(s)\approx\frac{V_{DD}}{I_0}
 \left(\frac{C_{\mathrm{ext}}}{s}+C_{\mathrm{int},0}\right).
 \label{eq:impl-sizing}
\end{equation}
در این مدل، \(C_{\mathrm{ext}}\) بار ثابت خارجی و \(C_{\mathrm{int},0}\) سهم داخلی
در اندازه اولیه است. بزرگ‌کردن گیت، سهم بار ثابت را سریع‌تر راه‌اندازی می‌کند، اما
سهم خازن داخلی نیز رشد می‌کند. خازن ورودی بزرگ‌تر، تأخیر و انرژی مرحله قبل را تغییر
می‌دهد. انتخاب اندازه به مسیر چندمرحله‌ای، \lr{Fan-in}، \lr{Fan-out}، سیم و محدودیت
مساحت وابسته است. مدل‌های دقیق‌تر جریان و خازن در فصل ۲ و روش‌های اندازه‌گذاری در فصل ۴ آمده‌اند.

\subsection{استخراج انرژی و توان از شارژ خازن}
توان لحظه‌ای منبع تغذیه برابر \(p=V_{DD}i\) است. برای شارژ کامل یک خازن از صفر
تا \(V_{DD}\)، منبع بار \(Q=C_LV_{DD}\) تحویل می‌دهد. چون ولتاژ منبع ثابت است،
انرژی دریافت‌شده از منبع برابر می‌شود با
\begin{equation}
 E_{\mathrm{supply}}=V_{DD}Q=C_LV_{DD}^2.
 \label{eq:impl-supply-energy}
\end{equation}
انرژی ذخیره‌شده در خود خازن از جمع انرژی‌های کوچک \(V\,dQ\) به دست می‌آید.
با \(dQ=C_L\,dV\)، مساحت زیر خط \(C_LV\) از صفر تا \(V_{DD}\) یک مثلث است، پس
\begin{equation}
 E_C=\int_0^{V_{DD}}C_LV\,dV=\frac12 C_LV_{DD}^2.
 \label{eq:impl-stored-energy}
\end{equation}
در شارژ معمول مقاومتی، نیم دیگر انرژی منبع در مسیر شارژ تلف می‌شود. هنگام تخلیه،
انرژی ذخیره‌شده نیز در مسیر تخلیه تلف می‌شود. بنابراین انرژی یک چرخه کامل
\(0\rightarrow1\rightarrow0\) در این مدل برابر \(C_LV_{DD}^2\) است.

اگر در هر ثانیه \(f_{\mathrm{cyc}}\) چرخه کامل رخ دهد، توان متوسط برابر انرژی
هر چرخه ضرب‌در تعداد چرخه در ثانیه خواهد بود
\begin{equation}
 P_{\mathrm{dyn}}=C_LV_{DD}^2f_{\mathrm{cyc}}.
 \label{eq:impl-dynamic-power}
\end{equation}
این همان صورت رابطه کتاب است \cite[صفحه ۴۳]{weste}. برای ارتباط با \lr{Clock}،
\(\alpha_{01}\) را متوسط تعداد گذارهای \(0\rightarrow1\) در هر دوره \lr{Clock}
تعریف می‌کنیم. در حالت پایدار، فرکانس شارژ \(\alpha_{01}f_{\mathrm{clk}}\) است، پس
\begin{equation}
 P_{\mathrm{dyn}}=\alpha_{01}C_LV_{DD}^2f_{\mathrm{clk}}.
 \label{eq:impl-activity-power}
\end{equation}
با شمارش مجموع گذارهای صعودی و نزولی، باید تعریف ضریب فعالیت را نیز عوض کرد.
برای سیگنال متناوب \(0,1,0,1,\ldots\) در نمونه‌های متوالی \lr{Clock}، در هر دو
دوره یک شارژ رخ می‌دهد، بنابراین \(\alpha_{01}=1/2\) است. برای چند گره، سهم
\(\alpha_{01,k}C_kV_{DD}^2f_{\mathrm{clk}}\) هر گره جمع می‌شود.

در مثال فرضی \(C_L=\qty{20}{\femto\farad}\)، \(V_{DD}=\qty{1}{\volt}\)،
\(f_{\mathrm{clk}}=\qty{100}{\mega\hertz}\) و \(\alpha_{01}=0.25\)، داریم
\begin{equation}
 P_{\mathrm{dyn}}=0.25(20\times10^{-15})(1)^2(100\times10^6)
 \,\unit{\watt}=\qty{0.50}{\micro\watt}.
 \label{eq:impl-power-example}
\end{equation}
کاهش ولتاژ به \qty{0.8}{\volt} با ثابت‌فرض‌کردن سایر عوامل، توان را در
\(0.8^2=0.64\) ضرب می‌کند و به \qty{0.32}{\micro\watt} می‌رساند. جریان ترانزیستور
نیز به ولتاژ وابسته است، پس بررسی سرعت به مدل جریان نیاز دارد.

\subsection{توان \lr{Static} و انتخاب پیاده‌سازی}
ترانزیستور خاموش واقعی جریان نشتی دارد. اگر جریان متوسط مستقیم از تغذیه به زمین
را \(I_{\mathrm{static}}\) بنامیم، از تعریف توان به دست می‌آید
\begin{equation}
 P_{\mathrm{static}}=V_{DD}I_{\mathrm{static}}.
 \label{eq:impl-static-power}
\end{equation}
این رابطه حتی در نبود تغییر منطقی خروجی کاربرد دارد \cite[صفحه ۴۳]{weste}.
برای جریان فرضی \qty{10}{\nano\ampere} و تغذیه \qty{1}{\volt}، توان
\qty{10}{\nano\watt} است. مجموع این سهم و مثال \eqref{eq:impl-power-example}،
در مدل دو مؤلفه‌ای، \qty{0.51}{\micro\watt} می‌شود. جریان گذرای هم‌پوشانی روشن‌شدن
شبکه‌های \lr{Pull-up} و \lr{Pull-down} نیز می‌تواند سهم دیگری ایجاد کند که در تحلیل
تفصیلی فصل توان بررسی می‌شود.

\lr{Synthesis} با انتخاب گیت‌های کتابخانه و اندازه آن‌ها، تابع منطقی را تحت
محدودیت‌های طراحی پیاده می‌کند. \lr{Custom Circuit Design} امکان انتخاب شبکه
ترانزیستوری متناسب با تابع را فراهم می‌سازد. کتاب اهمیت آن را برای بخش‌های دارای
نیاز سرعت یا توان ویژه و برای حجم تولید توجیه‌کننده هزینه مهندسی توضیح می‌دهد.
\lr{Area}، توان، سرعت و تلاش طراحی باید هم‌زمان سنجیده شوند.

\subsection{مثال \lr{Carry} با ۲۶ و ۱۲ ترانزیستور}
در \lr{Full Adder}، ورودی‌های تک‌بیتی \(a,b,c\) هستند و \(c\) نقش \lr{Carry-in}
را دارد. خروجی \(c_{\mathrm{out}}\) وقتی یک می‌شود که دست‌کم دو ورودی یک باشند.
تابع آن برابر است با
\begin{equation}
 c_{\mathrm{out}}=ab+ac+bc.
 \label{eq:impl-carry-majority}
\end{equation}
در این معادله، کنارهم‌نوشتن نمادها به معنای \lr{AND} و علامت جمع به معنای
\lr{OR} است. برای مثال \(a=b=1\) با هر مقدار \(c\)، خروجی یک می‌دهد.
اگر فقط یک ورودی یک باشد، هر سه حاصل‌ضرب صفر می‌شوند.

پیاده‌سازی شکل \secnum{۱.۶۱} کتاب از سه \lr{AND2} و یک \lr{OR3} تشکیل می‌شود
\cite{weste}. \lr{AND2} در \lr{Static CMOS} معمول، یک \lr{NAND2} چهارترانزیستوری
و یک \lr{Inverter} دوترانزیستوری دارد. \lr{OR3} نیز از \lr{NOR3} شش‌ترانزیستوری
و یک \lr{Inverter} ساخته می‌شود. بنابراین
\begin{equation*}
 N_{\mathrm{gate}}=3(4+2)+(6+2)=26.
\end{equation*}
مسیر ورودی تا خروجی از چهار مرحله وارون‌کننده منطقی عبور می‌کند.
تأخیر هر مرحله تابع اندازه و بار آن است.

برای ساخت مدار فشرده، ابتدا تابع را عامل‌گیری می‌کنیم
\begin{equation}
 F=ab+c(a+b),\qquad c_n=\overline F,\qquad c_{\mathrm{out}}=\overline{c_n}.
 \label{eq:impl-factored-carry}
\end{equation}
شبکه \lr{Pull-down} گره \(c_n\)، یک شاخه دوترانزیستوری سری با کنترل \(a,b\)
و یک شاخه شامل \(c\) در سری با دو ترانزیستور موازی \(a,b\) دارد. این شبکه پنج
\lr{nMOS} می‌خواهد. هرگاه \(F=1\) باشد، مسیری تا زمین برقرار می‌شود.

شبکه \lr{Pull-up} شکل کتاب نیز پنج \lr{pMOS} دارد. شرط هدایت آن را با
\(\bar a,\bar b,\bar c\) بیان می‌کنیم، زیرا \lr{pMOS} با ورودی صفر روشن می‌شود.
با قانون \lr{De Morgan} و توزیع حاصل‌ضرب داریم
\begin{align*}
 \overline F
 &=(\bar a+\bar b)(\bar c+\bar a\bar b)\\
 &=\bar c(\bar a+\bar b)+\bar a\bar b.
\end{align*}
جمله اول، \(c\) در سری با موازی \(a,b\)، و جمله دوم، سری \(a,b\) در شبکه
\lr{pMOS} را مشخص می‌کند. بنابراین برای هر ورودی دودویی پایدار، یکی از شبکه‌ها
گره \(c_n\) را به تغذیه مناسب می‌رساند. \lr{Inverter} پایانی دو ترانزیستور دیگر دارد
و تعداد کل \(5+5+2=12\) می‌شود. مسیر منطقی شامل یک گیت مرکب و یک
\lr{Inverter} است. مقایسه سرعت عددی، به اندازه ترانزیستورها و بارها نیاز دارد.

جدول \ref{tab:impl-carry} همه حالت‌های دودویی را کنترل می‌کند. ستون \lr{PDN}
و \lr{PUN}، وجود مسیر رسانا برای گره \(c_n\) را نشان می‌دهند.
\begin{table}[ht]
\centering
\caption{کنترل کامل تابع و شبکه‌های \lr{Carry} در مدل سوئیچی ایدئال.}
\label{tab:impl-carry}
\begin{tabular}{ccccccc}
\toprule
\(a\)&\(b\)&\(c\)&\lr{PDN}&\lr{PUN}&\(c_n\)&\(c_{\mathrm{out}}\)\\
\midrule
۰&۰&۰&۰&۱&۱&۰\\
۰&۰&۱&۰&۱&۱&۰\\
۰&۱&۰&۰&۱&۱&۰\\
۰&۱&۱&۱&۰&۰&۱\\
۱&۰&۰&۰&۱&۱&۰\\
۱&۰&۱&۱&۰&۰&۱\\
۱&۱&۰&۱&۰&۰&۱\\
۱&۱&۱&۱&۰&۰&۱\\
\bottomrule
\end{tabular}
\end{table}

\subsection{\lr{Schematic} و دو سطح \lr{Netlist}}
\lr{Schematic} اتصال‌ها را به صورت تصویر و \lr{Netlist} همان ارتباط‌ها را به صورت
متن نمایش می‌دهد. در توصیف ساختاری گیت‌ها، نمونه آموزشی تابع \lr{Carry} چنین است
\cite[صفحات ۴۳ و ۴۴]{weste}.
\begin{latin}
\small
\begin{verbatim}
module carry(input logic a, b, c, output logic cout);
  logic x, y, z;
  and g1(x, a, b);
  and g2(y, a, c);
  and g3(z, b, c);
  or  g4(cout, x, y, z);
endmodule
\end{verbatim}
\end{latin}
در این \lr{Verilog}، نخستین آرگومان هر گیت خروجی است. نام‌های \code{g1} تا
\code{g4} شناسه نمونه‌ها و \code{x}، \code{y} و \code{z} سیم‌های میانی هستند.
گیت عمومی هنوز اندازه، خازن و مدل فناوری را مشخص نمی‌کند.

در توصیف سوئیچی کتاب، \code{tranif1} با کنترل یک و \code{tranif0} با کنترل صفر
اتصال برقرار می‌کنند. ترتیب آرگومان آن‌ها دو پایانه مسیر و سپس کنترل است.
برای نمونه \code{tranif1 n3(cn, i1, c)} گره‌های \code{cn} و \code{i1} را هنگام
یک‌بودن \code{c} وصل می‌کند. مسیر در این مدل دوطرفه است. \code{supply0} و
\code{supply1} به‌ترتیب زمین و تغذیه منطقی، و \code{tri} گره سیمی با قابلیت حالت
پرامپدانس را معرفی می‌کنند. این نمایش برای بررسی اتصال و عملکرد سوئیچی مناسب است.

برای تحلیل پیوسته ولتاژ، جریان، تأخیر و توان، \lr{SPICE} به مدل افزاره و ابعاد نیاز
دارد. ترانزیستور \lr{MOS} چهار پایانه \lr{Drain}، \lr{Gate}، \lr{Source} و
\lr{Body} دارد. خط معرفی آن در کتاب به ترتیب زیر است
\cite[صفحات ۴۴ و ۴۵]{weste}.
\begin{latin}
\small
\begin{verbatim}
Mname drain gate source body model W=width L=length
+ AD=drain_area AS=source_area
Cname node1 node2 value
\end{verbatim}
\end{latin}
علامت \code{+} در آغاز خط، ادامه همان دستور را نشان می‌دهد. \code{W} و
\code{L} عرض و طول کانال، و \code{AD} و \code{AS} مساحت‌های اتصال
\lr{Drain} و \lr{Source} هستند. نام مدل، پارامترهای فناوری را به نمونه مرتبط می‌کند.
این ترتیب با ترتیب آرگومان گیت‌های \lr{Verilog} تفاوت دارد.

فهرست \lr{SPICE} آموزشی زیر اتصال و ابعاد مثال کتاب را حفظ می‌کند. پسوندهای
\code{u}، \code{p} و \code{f} به‌ترتیب ضرایب \(10^{-6}\)، \(10^{-12}\) و
\(10^{-15}\) هستند. بنابراین \code{W=2u} طول \qty{2}{\micro\metre} و
\code{AD=1.8p} مساحت \qty{1.8}{\micro\metre\squared} را مشخص می‌کنند.
نام‌های گره و ترتیب درگاه‌ها از مرجع حفظ شده‌اند.
\begin{latin}
\footnotesize
\begin{verbatim}
.SUBCKT CARRY A B C COUT VDD GND
MN1 I1 A GND GND NMOS W=2u L=0.6u AD=1.8p AS=3p
MN2 I1 B GND GND NMOS W=2u L=0.6u AD=1.8p AS=3p
MN3 CN C I1 GND NMOS W=2u L=0.6u AD=3p AS=3p
MN4 I2 B GND GND NMOS W=2u L=0.6u AD=0.9p AS=3p
MN5 CN A I2 GND NMOS W=2u L=0.6u AD=3p AS=0.9p
MP1 I3 A VDD VDD PMOS W=4u L=0.6u AD=3.6p AS=6p
MP2 I3 B VDD VDD PMOS W=4u L=0.6u AD=3.6p AS=6p
MP3 CN C I3 VDD PMOS W=4u L=0.6u AD=6p AS=6p
MP4 I4 B VDD VDD PMOS W=4u L=0.6u AD=1.8p AS=6p
MP5 CN A I4 VDD PMOS W=4u L=0.6u AD=6p AS=1.8p
MN6 COUT CN GND GND NMOS W=4u L=0.6u AD=6p AS=6p
MP6 COUT CN VDD VDD PMOS W=8u L=0.6u AD=12p AS=12p
CI1 I1 GND 6f
CI3 I3 GND 9f
CA A GND 12f
CB B GND 12f
CC C GND 6f
CCN CN GND 12f
CCOUT COUT GND 6f
.ENDS CARRY
\end{verbatim}
\end{latin}
در نمونه \code{MN3}، \lr{Body} به زمین متصل است، هرچند \lr{Source} روی گره
داخلی \code{I1} قرار دارد. \lr{Body}های \lr{pMOS} به \code{VDD} وصل شده‌اند.
مدل ترانزیستور، سهم خازن‌های ذاتی و پیوندی را با استفاده از ابعاد محاسبه می‌کند.
خازن‌های صریح این مثال، سهم اضافی سیم‌کشی را نمایش می‌دهند. در شبکه‌های دارای
سیم بلند، مقاومت سیم نیز باید در مدل مناسب منظور شود.

این قطعه، تعریف \lr{Subcircuit} است. اجرای تحلیل به تعریف معتبر مدل‌های
\code{NMOS} و \code{PMOS}، منبع تغذیه، تحریک ورودی، بار و دستور تحلیل نیاز دارد.
پارامترهای مدل فناوری باید از منبع فناوری انتخاب‌شده تأمین شوند.

\subsection{تمرین‌های پاسخ‌دار \lr{Circuit Design}}
\textbf{تمرین ۱.} بار \qty{30}{\femto\farad} با تغییر ولتاژ \qty{0.8}{\volt}
در \qty{120}{\pico\second} شارژ می‌شود. جریان متوسط لازم چقدر است؟
از \eqref{eq:impl-charge-delay} با حل برای جریان داریم
\[
 I_{\mathrm{avg}}=\frac{30\times10^{-15}\times0.8}{120\times10^{-12}}
 \,\unit{\ampere}=\qty{200}{\micro\ampere}.
\]
توان \(10^{-15}/10^{-12}=10^{-3}\) و ضریب \(24/120=0.2\) است، بنابراین
\(0.2\times10^{-3}\,\unit{\ampere}=\qty{200}{\micro\ampere}\) می‌شود.

\textbf{تمرین ۲.} اگر خازن دو برابر، فرکانس نصف و ولتاژ ۱۰ درصد کمتر شود، با
ضریب فعالیت ثابت، نسبت توان جدید به قبلی چقدر است؟ با تقسیم دو بار
\eqref{eq:impl-activity-power} بر یکدیگر داریم
\[
 \frac{P_2}{P_1}=\frac{C_2}{C_1}
 \left(\frac{V_2}{V_1}\right)^2\frac{f_2}{f_1}
 =2(0.9)^2\frac12=0.81.
\]
پس توان ۱۹ درصد کمتر می‌شود. این نتیجه به فرض‌های ثابت‌بودن ضریب فعالیت و
مستقل‌بررسی‌کردن پارامترهای داده‌شده وابسته است.

\textbf{تمرین ۳.} برای \(a=1,b=0,c=1\)، مسیر مدار فشرده را توضیح دهید.
در \lr{PDN}، شاخه \(c(a+b)\) از طریق ترانزیستورهای \(c\) و \(a\) وصل است.
در \lr{PUN}، ترانزیستورهای کنترل‌شده با \(a\) و \(c\) خاموش‌اند و مسیر کاملی
تا تغذیه وجود ندارد. پس \(c_n=0\) و \lr{Inverter} پایانی، \(c_{\mathrm{out}}=1\) می‌سازد.
\clearpage
